% Daftar Pustaka (gaya APA edisi ke-7) -- dibuat otomatis dari 00-common/pustaka.bib
\babtp{DAFTAR PUSTAKA}
Daftar pustaka berisi \textbf{semua} sumber yang dikutip di dalam teks, dan hanya sumber tersebut. Gunakan gaya APA edisi ke-7. Daftar ini \textbf{tidak diketik manual}: setiap sumber dicatat satu kali di berkas \texttt{00-common/pustaka.bib}, dipanggil di dalam teks dengan \verb|\parencite{kunci}| atau \verb|\textcite{kunci}|, lalu LaTeX menyusun daftar pustaka secara otomatis (diurutkan menurut abjad nama penulis pertama atau judul bila tanpa penulis, tanpa nomor urut; baris kedua dan seterusnya menjorok ke dalam). Mendeley atau Zotero dapat dipakai untuk mengekspor berkas \texttt{.bib}.

Contoh entri \texttt{pustaka.bib} mencakup jenis sumber yang umum dipakai: buku, buku terjemahan, bab dalam buku suntingan, artikel jurnal (dengan DOI), prosiding, skripsi atau tesis, laporan lembaga, lembar data (\textit{datasheet}), artikel daring, kode sumber di repositori, dan sumber tanpa penulis. Ganti dengan sumber Anda sendiri; sumber yang tidak dikutip tidak akan muncul di daftar ini.

\textbf{Sumber daring dicantumkan tanggal aksesnya.} Isi field \texttt{urldate} (format \texttt{2026-10-21}) untuk laporan daring, artikel daring, dan repositori; daftar pustaka akan menulis \textit{Diakses pada tanggal bulan tahun, dari alamat URL}, karena isinya dapat berubah. Tanggal akses adalah hari saat Anda membuka sumber tersebut. Artikel jurnal yang memiliki DOI cukup ditulis dengan DOI tanpa tanggal akses. Bila tahun terbit tidak diketahui, kosongkan field \texttt{date}; akan tertulis (t.t.).

Bentuk sitasi di dalam teks untuk jenis sumber di atas:
\begin{itemize}
\item satu penulis: \textcite{andrews2007} atau \parencite{andrews2007};
\item dua penulis: \parencite{xu2012};
\item tiga penulis atau lebih: \parencite{capozzi2013};
\item lembaga: \parencite{bps1997};
\item kutipan langsung dengan halaman: \parencite[15]{ghosh2010};
\item tanpa penulis: \parencite{panduan2020}.
\end{itemize}

% Sumber contoh yang belum dikutip di teks, agar tetap tampil. HAPUS baris \nocite ini pada naskah Anda.
\nocite{penulis2020,penulis2022,penulis2024,penulis2026,penulis2024kode,komdigi2025}

\begin{daftarpustaka}
\printbibliography[heading=none]
\end{daftarpustaka}
